\par\Needspace{6\baselineskip}\medskip
\begin{center}\hypertarget{native-245-article174}{}\pdfbookmark[2]{174.}{article-174}174.\end{center}
If the orbit of a celestial body has been determined according to the methods given in the preceding sections from three or four geocentric positions, each of which was derived from several observations as directed in the preceding article. that orbit will, as it were, occupy the middle ground among all these observations, and no trace of a pattern will remain in the differences between the observed and calculated places that could be removed or appreciably diminished by correcting the elements. Whenever the entire set of observations covers an interval of time that is not too long, we can in this way obtain the desired agreement of the elements with all the observations, provided the three or four so-called normal positions are judiciously chosen. In determining the orbits of comets or new planets whose observations do not yet extend beyond one year, this method will generally take us as far as the nature of the matter allows. Whenever, therefore, the orbit to be determined is inclined at a considerable angle to the ecliptic, it will generally be based on three observations chosen as far apart as possible: if, however, we happen to encounter in this way one of the cases excluded above (arts.\ 160---162), or if the inclination of the orbit appears too small, we shall prefer a determination from four positions, again taken as far apart as possible.

When a longer series of observations covering several years is already available, more normal positions can be derived from it: we should therefore make poor provision for the highest precision if we selected only three or four positions to determine the orbit and disregarded all the others entirely. On the contrary, if we aim at the utmost precision in such a case, we will endeavor to collect and use as many carefully determined positions as possible. We will then have more data than are required for determining the unknowns: but all these data will be subject to errors, however small, so that in general it is impossible to satisfy them all exactly. Now there is no reason to regard any particular six of these data as absolutely exact, and according to the principles of probability we must instead assume that greater or smaller errors are equally possible in \SourceLeaf{246}all of them without distinction; furthermore, speaking generally, slight errors occur more often than serious ones: clearly, then, an orbit that satisfies six data exactly but deviates more or less from the rest must be judged less consistent with the principles of the calculus of probabilities than one that differs somewhat from those six as well but agrees so much better with the rest. The investigation of the orbit with the \emph{greatest} probability in the strict sense will depend on knowing the law by which the probability of errors decreases as they grow larger: knowledge of this law, however, depends on so many vague or doubtful considerations --- physiological ones among them --- which cannot be subjected to calculation, that such a law can scarcely, if at all, be properly specified in any case in practical astronomy. Nevertheless, the investigation of the connection between this law and the most probable orbit, which we shall now undertake in full generality, is by no means to be regarded as idle speculation.
